% ch3 B.2 前半 (原书页 112–118 / PDF p127–p133)

\subsection{$N+1$ 体问题中声子的效应}

在以上全部讨论中，我们一直仿佛假定多体系统里只有价带中的电子可以自由运动。这一假设基于
Born-Oppenheimer 绝热近似；按照该近似，由于原子核与电子的质量相差极大，原子核的一切运动
相对电子而言都无限缓慢。对于为之发明的有限分子体系，这个近似非常好，因为其中电子激发能
从而其频率确实远高于离子的相应量，所以离子可以认为是静止的；但它当然恰恰会在我们觉得最
有意思的那个区域完全失效，即固体最低能量激发的区域，因为在这里我们无法\emph{先验地}断定
电子能量相对离子能量究竟是大还是小。

固体中在这里有任何重要性的离子运动，可以用离子实绕其平衡位置振动的量子化波动来描述
（图 32）。正如各位所知，这些量子化的波称为声子（phonon）。

%FIG{32}: 图 32

它们本身就是固体的元激发；事实上，它们是最早被发现的一组元激发——由 Einstein 和
Debye 在 1910 年前后发现。尽管如此，让你可能感到意外的是，把声子当作固体的多体集体激发
来处理时，仍有许多重要问题悬而未决。不过眼下我们只关心声子对我们一直在讨论的准粒子激发
性质的修正。

正如各位所知，对于每个 Bravais 原胞只含一个原子的固体，声子谱有三支，在长波长和简单
晶体取向下退化为一支纵波和两支横波声波；长波情形下频率由 $\omega_{k}=c_{l}k$ 或
$c_{t}k$ 给出，$c$ 为纵向或横向声速；能量为 $\hbar\omega_{k}$。

更复杂的固体的声子谱还有三支或更多的"光学"支，在长波长下退化为 Bravais 原胞内不同
原子彼此相对的运动；而在两个原子并不相同的情形下，

\begin{equation*}
(\text{Na})^{+}\!\to\;\; \leftarrow\!\!(\text{Cl})^{-}
\qquad\quad
(\text{Na})^{+}\!\to\;\; \leftarrow\!\!(\text{Cl})^{-}
\end{equation*}

单位体积便带有相应的偶极矩 $\mathbf{P}$。稍后我们会讨论到，在光学支中，纵波与横波的频率
在无限长波处相差一个有限量 $\omega_{lo}-\omega_{to}$。于是我们在图 33 中草绘了单原子
晶格和离子性双原子晶格的声子谱。

%FIG{33}: 图 33

声子的存在显然会带来若干后果。首先也是最简单的一点：我们看到，准粒子态不再必然是具有
给定动量或 $k$ 矢量的最低能量激发态；事实上，相对声子态而言，它通常将是一个能量非常高的
态。这正是我们在定义元激发时加以限定、要求激发态属于\emph{某一特定类型}的原因。就准粒子
与声子相对照的情形而言，区分其实相当容易，因为准粒子算符是费米型的，它使电子数目改变
1，即从 $N$ 变到 $N+1$，而产生声子的算符 $b_{\mathbf{K}}^{\dagger}$ 则近似是玻色型的，
既不改变离子数也不改变电子数；但一般说来区分并非如此简单。例如，态 $q_{k}^{\dagger}\Psi_{N}$
的能量未必低于带有一个声子加一个准粒子激发的态
$q_{k-\mathbf{K}}^{\dagger}b_{\mathbf{K}}^{\dagger}\Psi_{N}$；两个算符都是费米型的，
两个态都是 $N+1$ 粒子态。

当我们引入电子-声子相互作用的概念时，一个严重得多的问题出现了。这一相互作用有两种效应：
\emph{散射}与\emph{极化}。不过在讨论这些效应之前，我想先就电子-声子相互作用本身说几句话。

为了对这个相互作用可能取的形式获得一些概念，让我们在一个简单近似下把它计算出来。有一种
办法非常简单，但必须谨慎对待，这就是所谓的"刚性离子"（rigid-ion）模型。在此我们假定，
电子与之相互作用的晶格势是每个离子各自贡献之和：$V(r)=\sum_{j}V(r-\mathbf{R}_{j})$，或者
写成电子的二次量子化形式
$V=\int\mathrm{d}r\,\sum_{j}V(r-\mathbf{R}_{j})\,\Psi^{\dagger}(r)\Psi(r)$。我们为电子选用的
未微扰波函数，是在把 $\mathbf{R}_{j}$ 固定在格点上时算出的；于是微扰必定来自 $\mathbf{R}_{j}$
的位移 $\delta\mathbf{R}_{j}$，它由声子的存在引起：

\begin{equation*}
\mathcal{H}' \;=\; \sum_{j'} \int \mathrm{d}r\;\; \delta\mathbf{R}_{j} \cdot \nabla V(r-\mathbf{R}_{j})\, \Psi^{\dagger}(r)\, \Psi(r)
\end{equation*}

（译注：原书求和指标印作 $j'$，而式内各下标均为 $j$，两者应统一。）

$\delta\mathbf{R}_{j}$ 可以用声子振幅 $Q_{\mathbf{K}}$ 展开：

\begin{equation*}
\delta\mathbf{R}_{j} \;=\; \text{const} \cdot \sum_{\mathbf{K}} \mathrm{e}^{-i\mathbf{K}\cdot\mathbf{R}_{j}}\, Q_{\mathbf{K}}
\end{equation*}

而各个 $\Psi$ 则用 Bloch 函数展开：

\begin{align*}
\mathcal{H}' = {} & \int \mathrm{d}r \sum_{\substack{nn'k \\ k'\mathbf{K}j}}
\mathrm{e}^{-i\mathbf{K}\cdot\mathbf{R}_{j}}\,
\mathrm{e}^{ik\cdot r}\, u^{n*}_{k}(r)\,
\mathrm{e}^{-ik'\cdot r}\, u^{n'}_{k'}(r) \\
& \times\; Q_{\mathbf{K}}\, c^{n\dagger}_{k}\, c^{n'}_{k'} \cdot \nabla V(r-\mathbf{R}_{j}) \\
= {} & \sum_{k\mathbf{K}nn'} M^{nn'}_{k,\mathbf{K}}\, Q_{\mathbf{K}}\, c^{n'\dagger}_{k+\mathbf{K}}\, c^{n}_{k}
\end{align*}

其中

\begin{equation*}
M^{nn'}_{k\mathbf{K}} \;=\; \text{const.} \int \nabla V(r-\mathbf{R}_{j})\, \phi^{n}_{k+\mathbf{K}}(r)\, \phi^{n'}_{k}(r)\;\, \mathrm{d}\Omega_{\mathrm{cell}}
\tag{7}
\end{equation*}

$Q_{\mathbf{K}}$ 又可以用声子的产生与湮灭算符 $b_{\mathbf{K}}$ 和 $b_{-\mathbf{K}}^{\dagger}$
表出。这样一来，相互作用就具有电子被散射同时产生或湮灭一个声子的形式（图 34）。

%FIG{34}: 图 34

比这个相当朴素的模型更为认真的分析其实是可以做出的；关于这些分析的精彩论述，我向各位
推荐 Ziman 的书（文献 (3)）中相应的章（第五章）。这里我只想集中讨论一个极限情形，即
$k$ 很小的情形。

查阅 Ziman 的书可以发现，在这个极限下矩阵元 $M$ 结果正比于 $\mathbf{K}$，于是相互作用
并不像乍看之下那样正比于离子的位移 $\delta\mathbf{R}_{j}\propto Q_{\mathbf{K}}$，而是正比于
其应变（或称"形变"）$\nabla\mathbf{R}_{j}\propto\mathbf{K}Q_{\mathbf{K}}$（译注：原文如此，
按本意应为 $\nabla\,\delta\mathbf{R}_{j}$）。在单电子理论范围内证明这一点并不十分困难，但
有趣的是要证明它在一般情形下也成立，即使把多体效应考虑进去。

这非常值得去做，因为脱离价带电子来谈论离子和离子实自身的位移是没有意义的——至少在某种
意义上，这些电子是随它们一起被携带的。撇开别的不论，我们至少意识到：由于离子实带有净电荷，
当它们形成长波长的纵波时，会在波峰处积累起非常大的空间电荷，这种电荷必然要移动价电子的
电荷分布。一般地说，"价电子气整体不随离子运动"这一假设显然站不住脚，这一点是明显的。
因此，电子-声子相互作用中有相当大的部分必须被认为已经包含在声子激发本身的定义之中，正如
被散射的电子其实并不是单个的独立电子，而是准粒子一样。

这个问题与场论中的重整化问题极为相似；我们看到的"物理"准粒子和声子，与我们能够简单
设想的"裸"粒子全然不同，其实验性质——能量、相互作用等等——包含了围绕裸粒子的扰动云的
贡献。在固体理论中，我们并没有场论中那种不可挽回地阻止我们找出"裸"性质与"物理"性质之间
真实关系的发散；尽管如此，固体是如此复杂的对象，以致为了求出的不仅有声子频率和准粒子
能量、还有它们之间的相互作用，求诸实验往往比求诸理论更有用。这就是\emph{形变势}
（deformation potential）理论（文献 (70)）背后的思想。在这一理论中，我们借助均匀应变的
可观察（至少原则上可观察）效应来估计长波长声子的效应。

于是设想，为简单起见，我们有一块绝缘体，并以某种相当平滑的方式使它发生弹性形变，于是，
若各处的局域位移为 $\delta\mathbf{R}_{j}$，则任一局域区域内都有明确定义的局域压缩
$\Delta=\nabla\cdot\delta\mathbf{R}_{j}$，并且在愿意的话还有切应变
$S=\operatorname{symm}\bigl\{\nabla\,\delta\mathbf{R}_{j}-\tfrac{1}{3}\,\nabla\cdot(\delta\mathbf{R}_{j})\,\mathbf{1}\bigr\}$。
我们让 $N$ 体电子气完全绝热地随之形变。于是，在新的、已形变的晶格中，我们可以构建一个
波包 $\Psi_{\text{packet}}(k_{0},r_{0})$，其中心动量为 $k_{0}$，位于围绕 $r_{0}$、尺度为
$1/\Delta k$ 的区域内。假定该区域内的形变 $\Delta$、$S$ 相当均匀；否则，晶格当然就不再
足够周期，我们也就无法在不确定度 $\Delta k$ 之内定义波矢 $k_{0}$。这个由精确准粒子构成的
新波包具有新的能量 $E^{\prime}_{n}(k_{0})$，它与原来的能量不同，因为晶格在局域上发生了
形变；若无形变，显然无论原子的\emph{位移}有多大，能量都根本不会改变。由于声子中原子的
运动相当小，可以合理地预期带能量对形变是线性的，其线性比例常数称为"形变势"常数 $E$：

\begin{equation*}
E^{\prime}_{n}(k_{0}) \;-\; E_{n}(k_{0}) \;=\; E_{\Delta}(k_{0})\,\Delta \;+\; E_{S}(k_{0})\, S \;+\; \cdots
\tag{8}
\end{equation*}

这一\emph{以经验方式定义的}势，就是准粒子 $q_{k}^{\dagger}$ 的真实散射势，它既包括形变给
$\Psi_{g}$ 带来的近乎无穷的改变，也包括准粒子波函数本身的修正。这可以按如下方式证明。
围绕动量 $k_{0}$ 和位置 $r$ 产生一个波包的算符是

\begin{equation*}
q_{k_{0}}^{\dagger}(r) \;=\; \int \mathrm{d}k\;\, f(k-k_{0})\; q_{k}^{\dagger}\; \mathrm{e}^{ik\cdot r}
\end{equation*}

而检验这样的波包是否已被产生的数算符是

\begin{equation*}
n_{k_{0}}(r) \;=\; q_{k_{0}}^{\dagger}(r)\, q_{k_{0}}(r)
\end{equation*}

于是，对一个实际波包波函数能产生与式 (8) 的能量相同效果的算符，可以取为（我们稍后将验证
它确实具有正确的效果）：

\begin{equation*}
\sum_{k_{0}} \int \mathrm{d}r\;\; E_{\Delta}(k_{0})\;\, \Delta(r)\, n_{k_{0}}(r) \;+\; \text{（对 } S \text{ 同样处理）}
\end{equation*}

把它算出来，得到

\begin{align*}
& \sum_{k_{0}\mathbf{K}kk'} \int \mathrm{d}r\;\, f(k-k_{0})\, f(k'-k_{0})\;
\mathrm{e}^{i(k-k')\cdot r}\; q_{k}^{\dagger} q_{k'},\;\, E_{\Delta}(k_{0})\;\, \mathrm{e}^{i\mathbf{K}\cdot r}\;\; \Delta_{\mathbf{K}} \\
& \qquad = \sum_{k_{0}k\mathbf{K}} f(k-k_{0})\, f(k+\mathbf{K}-k_{0})\;\, q_{k}^{\dagger} q_{k+\mathbf{K}}\;\; \Delta_{\mathbf{K}}\, E_{\Delta}(k_{0})
\end{align*}

我们必须假定 $\mathbf{K}$ 小于波包的展宽，否则各个 $f$ 因子将不能正确地交叠；这就是前面
已经表述过的要求，即波包必须限制在足够小的体积内，使 $\Delta$ 在其内部均匀。如果
$\mathbf{K}$ 足够小，在 $f(k+\mathbf{K}-k_{0})$ 中便可将它忽略，于是可以利用波包函数的
完备性关系

\begin{equation*}
\sum_{k_{0}} f^{2}(k-k_{0}) \;=\; 1
\end{equation*}

若波包展宽又比 $E_{\Delta}(k_{0})$ 随 $k_{0}$ 变化的范围小，我们最终得到

\begin{equation*}
(8) \;=\; \sum_{k\mathbf{K}} E_{\Delta}(k)\, q_{k}^{\dagger} q_{k+\mathbf{K}}\;\; \Delta_{\mathbf{K}}
\tag{9}
\end{equation*}

把这一表达式作用于相应的波包，即可予以验证：

\begin{align*}
& \sum_{k\mathbf{K}} E_{\Delta}(k)\, q_{k}^{\dagger} q_{k+\mathbf{K}}\;\; \Delta_{\mathbf{K}} \sum_{k'} f(k'-k_{0})\, q_{k'}^{\dagger}\, \mathrm{e}^{ik'\cdot r}\, \Psi_{g} \\
& \qquad = \sum_{k\mathbf{K}} E_{\Delta}(k)\, \Delta_{\mathbf{K}}\;\, \mathrm{e}^{i\mathbf{K}\cdot r}\, f(k+\mathbf{K}-k_{0})\, q_{k}^{\dagger}\, \mathrm{e}^{ik\cdot r}\, \Psi_{g}
\end{align*}

再者，如果 $\mathbf{K}\ll\Delta k$，即 $f$ 的宽度，那么上式就正好化为